\documentclass[11pt]{amsart}

\usepackage[T1]{fontenc}
\usepackage{lmodern}
\usepackage{microtype}
\usepackage{amsmath,amssymb,amsthm}
\usepackage[hidelinks]{hyperref}

\hypersetup{
  pdftitle={A Proof of Calderon's mod-p Bailey Congruence at Every Odd Inert Prime}
}

\newtheorem{theorem}{Theorem}
\newtheorem{lemma}[theorem]{Lemma}
\newtheorem{corollary}[theorem]{Corollary}
\theoremstyle{remark}
\newtheorem*{remark}{Remark}

\newcommand{\OO}{\mathcal{O}}
\newcommand{\Rw}{R_{\omega}}
\newcommand{\Fp}{\mathbb{F}_p}
\newcommand{\ZZ}{\mathbb{Z}}
\newcommand{\QQ}{\mathbb{Q}}
\newcommand{\legendre}[2]{\left(\frac{#1}{#2}\right)}

\title[Calderon's Bailey congruence at every odd inert prime]
{A Proof of Calderon's mod-$p$ Bailey Congruence at Every Odd Inert Prime}

\date{}

\subjclass[2020]{Primary 11B65; Secondary 11A07, 11R11, 05A10}
\keywords{Bailey congruence, Lucas' theorem, rectangular Gaussian binomial
coefficient, quadratic order, inert prime, Frobenius}

\begin{document}

\begin{abstract}
Let $\omega$ be a root of $x^{2}-Tx+N$ with $T,N\in\ZZ$, let
$\OO=\ZZ[\omega]$, and for $A\ge C\ge1$, $B\ge D\ge1$ let
$\Rw(A,B;C,D)$ be the rectangular coefficient of Calderon
[arXiv:2608.00347v1]. His Conjecture~6.5 asserts a mod-$p$ digit
factorisation of $\Rw$ for every imaginary quadratic order and every prime
satisfying his hypothesis~(6.3). We prove it, under the weaker hypothesis
that $p$ is odd, $p\nmid\Delta:=T^{2}-4N$ and $\legendre{\Delta}{p}=-1$;
neither $p>5$ nor $\Delta<0$ is used. The proof is elementary: the only
inputs are that $\{1,\omega\}$ is an $\Fp$-basis of
$\OO/p\OO\cong\mathbb{F}_{p^{2}}$ and that $\omega^{p}\equiv\bar\omega$
there. The mechanism is the one the author sketches in the twenty-five
lines above his conjecture; what this note supplies is the verification he
defers.
\end{abstract}

\maketitle

\section{The conjecture}

Fix $T,N\in\ZZ$ and let $\omega$ be a root of $x^{2}-Tx+N$ with conjugate
$\bar\omega=T-\omega$ and discriminant $\Delta=T^{2}-4N$. Put
$\OO=\ZZ[\omega]$, and for a finite box
$\mathcal B\subset\ZZ_{\ge1}\times\ZZ_{\ge1}$ put
\[
 F(\mathcal B)=\prod_{(u,v)\in\mathcal B}(u+v\omega)\in\OO .
\]
For $A\ge C\ge1$ and $B\ge D\ge1$ the \emph{rectangular coefficient} is
\begin{equation}\label{eq:R}
 \Rw(A,B;C,D)
 =\frac{F\bigl([A-C+1,A]\times[B-D+1,B]\bigr)}
        {F\bigl([1,C]\times[1,D]\bigr)}
 \in\QQ(\omega).
\end{equation}
For a prime $p$ write $\OO_{\omega,p}=\OO\otimes_{\ZZ}\ZZ_p
=\ZZ_p+\ZZ_p\omega$; note that this last equality holds for
\emph{every} $p$, because $\{1,\omega\}$ is a $\ZZ$-basis of $\OO$.

Calderon \cite{Calderon} states the following. We keep the author's
hypotheses and index names and set the display in our own notation.

\begin{quote}
\textbf{Conjecture 6.5.} Let $p$ satisfy (6.3), and suppose
$1\le\gamma\le\alpha\le p-1$ and $1\le\delta\le\beta\le p-1$. Then
\begin{equation}\label{eq:C65}
\begin{split}
 \Rw(pA+\alpha,&\,pB+\beta;pC+\gamma,pD+\delta)\\
 &\equiv
 \Rw(A,B;C,D)\,\Rw(\alpha,\beta;\gamma,\delta)
 \binom{\alpha}{\gamma}^{\!D}\binom{\beta}{\delta}^{\!C}
 \pmod{p\,\OO_{\omega,p}}.
\end{split}
\end{equation}
\end{quote}

Three clauses are inherited from the surrounding text rather than restated
in that environment. The defining relation of $\omega$ asks only for
$\omega^{2}-T\omega+N=0$ with $T,N\in\ZZ$ and $\Delta<0$; the ring $\OO_{\omega}:=\ZZ[\omega]$ is called
an order in an imaginary quadratic field, so \emph{maximality is never
assumed}. Hypothesis~(6.3) is
\begin{equation}\label{eq:63}
 p>5,\qquad p\nmid\Delta,\qquad \legendre{\Delta}{p}=-1 .
\end{equation}
The index ranges $A\ge C\ge1$, $B\ge D\ge1$ come from the definition
\eqref{eq:R}.

What \cite{Calderon} \emph{proves} is his Theorem~1.3, and only for
$\omega=i$, since the coefficient there is defined over $\ZZ[i]$ alone;
its hypothesis is that $p\equiv3\bmod4$ be an odd prime. For $\omega=i$
one has $\Delta=-4$ and $\legendre{-4}{p}=-1$ exactly when
$p\equiv3\bmod4$, so Theorem~1.3 is precisely the case $T=0$, $N=1$ of
what follows.

\begin{theorem}\label{thm:main}
Let $p$ be an odd prime and $T,N\in\ZZ$ with
\begin{equation}\label{eq:I}
 p\nmid\Delta
 \qquad\text{and}\qquad
 \legendre{\Delta}{p}=-1 .
\end{equation}
Let $A\ge C\ge1$, $B\ge D\ge1$, and let
$1\le\gamma\le\alpha\le p-1$, $1\le\delta\le\beta\le p-1$. Then all three
rectangular coefficients occurring in \eqref{eq:C65} lie in
$\OO_{\omega,p}$, and \eqref{eq:C65} holds.
In particular Conjecture~6.5 of \cite{Calderon} is true, and it remains
true when \eqref{eq:63} is weakened to \eqref{eq:I}.
\end{theorem}

\section{The proof}\label{sec:proof}

Assume \eqref{eq:I}. Then $x^{2}-Tx+N$ is irreducible over $\Fp$, so
$\mathfrak p:=p\,\OO_{\omega,p}$ is a maximal ideal, $\OO_{\omega,p}$ is a
discrete valuation ring with uniformiser $p$ and residue field
\[
 \OO_{\omega,p}/\mathfrak p\cong\mathbb F_{p^{2}},
\]
in which $\{1,\omega\}$ is an $\Fp$-basis. Write $v$ for the valuation
with $v(p)=1$, extended to $\QQ(\omega)^{\times}$, and write $\equiv$ for
congruence modulo $\mathfrak p$. We use exactly three facts.

\begin{itemize}
\item[(i)] For $a,b\in\ZZ$ not both zero,
  $v(a+b\omega)=\min\bigl(v_p(a),v_p(b)\bigr)$.

  Indeed, with $m=\min(v_p(a),v_p(b))$ write $a+b\omega=p^{m}(u'+v'\omega)$;
  then $u',v'$ are not both divisible by $p$, so $u'+v'\omega\not\equiv0$
  because $\{1,\omega\}$ is an $\Fp$-basis of the residue field, i.e.
  $u'+v'\omega$ is a unit of $\OO_{\omega,p}$.

\item[(ii)] $\omega^{p}\equiv\bar\omega$. The map $x\mapsto x^{p}$ is the
  nontrivial element of $\operatorname{Gal}(\mathbb F_{p^{2}}/\Fp)$, so it
  interchanges the two roots of $x^{2}-Tx+N$.

\item[(iii)] For $x,y\in\Fp$,
  \[
   \prod_{t\in\Fp}(x+t\omega)=c_{1}x,
   \qquad
   \prod_{t\in\Fp}(t+y\omega)=c_{2}y,
   \qquad
   c_{1}=\frac{\omega-\bar\omega}{\omega},\quad c_{2}=\bar\omega-\omega ,
  \]
  and $c_{1},c_{2}\ne0$. For the first, $\prod_{t\in\Fp}(X+t)=X^{p}-X$
  gives, on putting $X=x/\omega$ and multiplying by $\omega^{p}$,
  \[
   \prod_{t\in\Fp}(x+t\omega)=x^{p}-\omega^{p-1}x
   =x-\frac{\bar\omega}{\omega}x=c_{1}x
  \]
  by (ii) and $x^{p}=x$; the second is
  $\prod_{t}(t+y\omega)=(y\omega)^{p}-y\omega=y(\bar\omega-\omega)$.
  Finally $(\omega-\bar\omega)^{2}=\Delta\ne0$ in $\Fp$ by \eqref{eq:I},
  and $\omega$ is a unit by (i), so $c_{1}$ and $c_{2}$ are nonzero.
  Note that neither constant is evaluated below.
\end{itemize}

\begin{lemma}[integrality]\label{lem:int}
For all $A\ge C\ge1$ and $B\ge D\ge1$ we have
$v\bigl(\Rw(A,B;C,D)\bigr)\ge0$; that is, $\Rw(A,B;C,D)\in\OO_{\omega,p}$.
\end{lemma}

\begin{proof}
By (i) and $\min(a,b)=\#\{j\ge1:p^{j}\mid\text{both}\}$,
\[
 v\bigl(F(\mathcal I\times\mathcal J)\bigr)
 =\sum_{j\ge1}M_{\mathcal I}(p^{j})\,M_{\mathcal J}(p^{j}),
\]
where $M_{\mathcal I}(q)$ is the number of multiples of $q$ in the
interval $\mathcal I$. An interval of length $C$ contains at least
$\lfloor C/q\rfloor$ multiples of $q$, and $[1,C]$ contains exactly
$\lfloor C/q\rfloor$. Applying this to $\mathcal I=[A-C+1,A]$ and
$\mathcal J=[B-D+1,B]$ gives
$v\bigl(F([A-C+1,A]\times[B-D+1,B])\bigr)\ge
 v\bigl(F([1,C]\times[1,D])\bigr)$.
\end{proof}

\begin{proof}[Proof of Theorem \ref{thm:main}]
Write $A'=pA+\alpha$, $B'=pB+\beta$, $C'=pC+\gamma$, $D'=pD+\delta$, and
put $\mathcal I=[A'-C'+1,A']$, $\mathcal J=[B'-D'+1,B']$, so that
$|\mathcal I|=C'$ and $|\mathcal J|=D'$ and the left side of
\eqref{eq:C65} is
$F(\mathcal I\times\mathcal J)/F([1,C']\times[1,D'])$.

\smallskip
\noindent\emph{Step 1: split off the $p$-divisible factors, exactly.}
By (i), $p\mid(u+v\omega)$ if and only if $p\mid u$ and $p\mid v$. An
integer $pk$ lies in $\mathcal I$ if and only if
$p(A-C)+(\alpha-\gamma+1)\le pk\le pA+\alpha$; since
$1\le\alpha-\gamma+1\le p-1$ and $0\le\alpha\le p-1$, this says exactly
$A-C+1\le k\le A$, so $\mathcal I$ contains exactly the $C$ multiples
$\{pk:A-C+1\le k\le A\}$. Likewise $\mathcal J$ contains exactly
$\{pl:B-D+1\le l\le B\}$, and, because $1\le\gamma\le p-1$ and
$1\le\delta\le p-1$, the interval $[1,C']$ contains exactly
$\{pk:1\le k\le C\}$ and $[1,D']$ exactly $\{pl:1\le l\le D\}$. Therefore
\begin{align*}
 F(\mathcal I\times\mathcal J)
  &=p^{CD}\,F\bigl([A-C+1,A]\times[B-D+1,B]\bigr)\cdot \mathrm{Num},\\
 F\bigl([1,C']\times[1,D']\bigr)
  &=p^{CD}\,F\bigl([1,C]\times[1,D]\bigr)\cdot \mathrm{Den},
\end{align*}
where $\mathrm{Num}$ and $\mathrm{Den}$ are the products over those pairs
$(u,v)$ of the respective boxes that do \emph{not} have both coordinates
divisible by $p$. The powers $p^{CD}$ cancel, so as an identity in
$\QQ(\omega)^{\times}$,
\begin{equation}\label{eq:step1}
 \Rw(A',B';C',D')=\Rw(A,B;C,D)\cdot\frac{\mathrm{Num}}{\mathrm{Den}},
\end{equation}
and by (i) both $\mathrm{Num}$ and $\mathrm{Den}$ are units of
$\OO_{\omega,p}$. In particular
\[
 v\bigl(\Rw(A',B';C',D')\bigr)=v\bigl(\Rw(A,B;C,D)\bigr):
\]
the $p$-power bookkeeping cancels uniformly, on \eqref{eq:I} alone.

\smallskip
\noindent\emph{Step 2: a residue census.}
Modulo $\mathfrak p$ the factor $u+v\omega$ depends only on
$(u\bmod p,\ v\bmod p)$, so
\[
 \mathrm{Num}\equiv\prod_{(r,s)\ne(0,0)}(r+s\omega)^{e(r)f(s)},
 \qquad
 \mathrm{Den}\equiv\prod_{(r,s)\ne(0,0)}(r+s\omega)^{e_{0}(r)f_{0}(s)},
\]
the products running over $(r,s)\in\Fp\times\Fp$, where $e(r)$ is the
number of $u\in\mathcal I$ with $u\equiv r$ and $e_{0}$, $f$, $f_{0}$ are
defined likewise for $[1,C']$, $\mathcal J$, $[1,D']$. Since
$|\mathcal I|=pC+\gamma$ and $\mathcal I$ ends at $A'\equiv\alpha$,
\[
 e(r)=C+a_{r},\quad a_{r}=[\,r\in S_{\alpha}\,],\qquad
 S_{\alpha}:=\{\alpha-\gamma+1,\dots,\alpha\},
\]
and similarly $e_{0}(r)=C+a'_{r}$ with $a'_{r}=[\,r\in\{1,\dots,\gamma\}\,]$,
$f(s)=D+b_{s}$ with $b_{s}=[\,s\in S_{\beta}\,]$,
$S_{\beta}:=\{\beta-\delta+1,\dots,\beta\}$, and $f_{0}(s)=D+b'_{s}$ with
$b'_{s}=[\,s\in\{1,\dots,\delta\}\,]$. \emph{The residue $0$ lies in none
of the four sets $S_{\alpha}$, $\{1,\dots,\gamma\}$, $S_{\beta}$,
$\{1,\dots,\delta\}$} --- this is exactly what
$1\le\gamma\le\alpha\le p-1$ and $1\le\delta\le\beta\le p-1$ buy --- so
$a_{0}=a'_{0}=b_{0}=b'_{0}=0$, and the exponent of $(r+s\omega)$ in
$\mathrm{Num}/\mathrm{Den}$ is
\[
 e(r)f(s)-e_{0}(r)f_{0}(s)
 =C(b_{s}-b'_{s})+D(a_{r}-a'_{r})+(a_{r}b_{s}-a'_{r}b'_{s}).
\]

\smallskip
\noindent\emph{Step 3: three groups.}
Summing the three terms separately, and using $a_{0}=b_{0}=0$ to see that
the first two extend over all $s$ resp.\ all $r$,
\[
 \frac{\mathrm{Num}}{\mathrm{Den}}
 \equiv
 \Biggl[\frac{\prod_{r\in S_{\alpha}}V(r)}{\prod_{r=1}^{\gamma}V(r)}\Biggr]^{D}
 \cdot
 \Biggl[\frac{\prod_{s\in S_{\beta}}W(s)}{\prod_{s=1}^{\delta}W(s)}\Biggr]^{C}
 \cdot
 \frac{F(S_{\alpha}\times S_{\beta})}{F([1,\gamma]\times[1,\delta])},
\]
where $V(r)=\prod_{s\in\Fp}(r+s\omega)$ and
$W(s)=\prod_{r\in\Fp}(r+s\omega)$.

\smallskip
\noindent\emph{Step 4: cardinality does the rest.}
The third group is $\Rw(\alpha,\beta;\gamma,\delta)$ by the definition
\eqref{eq:R}, since $S_{\alpha}=[\alpha-\gamma+1,\alpha]$ and
$S_{\beta}=[\beta-\delta+1,\beta]$; it is a unit of $\OO_{\omega,p}$
because every entry of both boxes lies in $[1,p-1]$. By (iii),
$V(r)=c_{1}r$ and $W(s)=c_{2}s$, and $|S_{\alpha}|=\gamma=|\{1,\dots,\gamma\}|$,
so $c_{1}^{\gamma}$ cancels \emph{by cardinality alone}, leaving
\[
 \frac{\prod_{r\in S_{\alpha}}V(r)}{\prod_{r=1}^{\gamma}V(r)}
 =\frac{\alpha(\alpha-1)\cdots(\alpha-\gamma+1)}{\gamma!}
 =\binom{\alpha}{\gamma}
 \quad\text{in }\Fp,
\]
and symmetrically the second group contributes $\binom{\beta}{\delta}^{C}$.
Hence, writing $U=\mathrm{Num}/\mathrm{Den}$ and
$Z=\Rw(\alpha,\beta;\gamma,\delta)\binom{\alpha}{\gamma}^{D}
\binom{\beta}{\delta}^{C}$, we have $U\equiv Z$ with $U,Z$ units, and by
\eqref{eq:step1}
\[
 \Rw(A',B';C',D')-\Rw(A,B;C,D)\,Z
 =\Rw(A,B;C,D)\,(U-Z).
\]
Here $v(U-Z)\ge1$ and $v\bigl(\Rw(A,B;C,D)\bigr)\ge0$ by
Lemma~\ref{lem:int}, so the left side has valuation at least $1$, which is
\eqref{eq:C65}. Nothing in the argument used $\Delta<0$, $p>5$, the unit
group of $\OO$, or the maximality of $\OO$.
\end{proof}

\begin{remark}
Two of the four sets in Step~2 exclude $0$ only because
$\gamma\le\alpha\le p-1$. Drop the digit bounds and $0$ enters
$S_{\alpha}$, at which point $V(0)=0$ and Step~4 collapses.
\end{remark}

\begin{remark}
At a split $p$ fact (i) fails --- $p$ divides $u+v\omega$ for pairs with
$p\nmid u$ --- and the congruence fails in general: for $\omega=i$ and
$p=13$ it already fails at $(A,B,C,D)=(1,2,1,1)$ with
$\alpha=\beta=\gamma=\delta=1$. Theorem~\ref{thm:main} does \emph{not}
answer Problem~6.7 of \cite{Calderon}, which asks for a \emph{pair} of
local congruences at the two primes above a split $p$.
\end{remark}

\medskip
\noindent\textbf{Artifacts.}
Theorem~\ref{thm:main} is a hand proof, and the program accompanying this
note does not stand in for it. That program, \texttt{verify.py}, re-derives
the finite combinatorial content of Steps~1, 2 and 4 and of
Lemma~\ref{lem:int} over stated finite ranges; checks facts (i)--(iii) in
each family it uses; exhausts \eqref{eq:C65} itself, without failure, over
the index ranges it prints, in families satisfying \eqref{eq:I} both
inside and outside \eqref{eq:63}; and reproduces the split-prime failure
quoted in the remark above. It uses the Python standard library and exact
integer arithmetic alone, and reports what it does not cover. A recorded
run is included as \texttt{verify.output.txt}.

\begin{thebibliography}{9}

\bibitem{Calderon}
K.~Calderon,
\emph{Ljunggren--Jacobsthal and Bailey-type congruences for rectangular
Gaussian binomial coefficients},
arXiv:2608.00347v1, 31 July 2026.
\href{https://arxiv.org/abs/2608.00347}{arXiv:2608.00347}.

\bibitem{Bailey}
D.~F. Bailey,
\emph{Two $p^{3}$ variations of Lucas' theorem},
J. Number Theory \textbf{35} (1990), no.~2, 208--215.

\bibitem{Lucas}
E.~Lucas,
\emph{Th\'eorie des fonctions num\'eriques simplement p\'eriodiques},
Amer. J. Math. \textbf{1} (1878), 184--196, 197--240, 289--321.

\bibitem{Kalinin}
N.~Kalinin,
\emph{Wolstenholme's theorem over Gaussian integers},
Funct. Approx. Comment. Math. \textbf{74} (2026), no.~2, 213--223;
\href{https://doi.org/10.7169/facm/250713-21-7}
{doi:10.7169/facm/250713-21-7}.

\end{thebibliography}

\end{document}
